\documentclass[oneside,a4paper,14pt]{extarticle}
\usepackage[a4paper,letterpaper,top=20mm,bottom=20mm,left=20mm,right=10mm]{geometry}
\usepackage[russian]{babel}
\usepackage{indentfirst}
\usepackage{graphicx}
\usepackage{caption}
\usepackage{titlesec}
\usepackage{minted, fancyvrb}
\usepackage{hyperref}
\usepackage{enumitem}
\usepackage{amsmath}
\usepackage{fancyhdr}
\usepackage{colortbl}
\usepackage[dvipsnames]{xcolor}
\usepackage{tikz}
\usepackage{float}
\usepackage{longtable}

% Настройка автоматических подписей к рисункам
\DeclareCaptionLabelSeparator{customdash}{~---~}
\captionsetup[figure]{name=Рисунок, labelsep=customdash, justification=centering}

% Форматирование заголовков
\titleformat{\section}{\normalsize\bfseries}{\thesection}{1em}{}
\titleformat{\subsection}{\normalsize\bfseries}{\thesubsection}{1em}{}

% Интерлиньяж и абзац
\renewcommand\baselinestretch{1.33}
\setlength{\parindent}{1.25cm}

% Гиперссылки
\hypersetup{
  colorlinks=true,
  linkcolor=black,
  urlcolor=blue,
  pdfborder={0 0 0},
  pdftitle={Синтез 4-разрядного компаратора},
  pdfauthor={Черкасов А.А.},
}

\begin{document}

\newpage
\thispagestyle{empty}
\begin{center}
  МИНИСТЕРСТВО НАУКИ И ВЫСШЕГО ОБРАЗОВАНИЯ РОССИЙСКОЙ ФЕДЕРАЦИИ\\
  ФЕДЕРАЛЬНОЕ ГОСУДАРСТВЕННОЕ БЮДЖЕТНОЕ УЧРЕЖДЕНИЕ ВЫСШЕГО ОБРАЗОВАНИЯ\\
  <<ВЯТСКИЙ ГОСУДАРСТВЕННЫЙ УНИВЕРСИТЕТ>>\\
  Институт математики и информационных систем\\
  Факультет автоматики и вычислительной техники\\
  Кафедра электронных вычислительных машин
\end{center}
\vspace{10mm}

\hfill
\begin{tabular}{l}
  \footnotesize Дата сдачи на проверку:                                          \\
  \footnotesize <<\rule[-1mm]{5mm}{0.10mm}\/>>\rule[-1mm]{20mm}{0.10mm}\ 2026 г. \\
  \footnotesize Проверено:                                                       \\
  \footnotesize <<\rule[-1mm]{5mm}{0.10mm}\/>>\rule[-1mm]{20mm}{0.10mm}\ 2026 г. \\
\end{tabular}
\vfill

\begin{center}
  Отчёт по контрольной работе\\
  <<Синтез 4-разрядного цифрового компаратора>>\\
  по дисциплине\\
  <<Цифровая схемотехника>>\\
\end{center}
\vspace{25mm}
\noindent
\begin{tabular}{ll}
  Разработал студент гр. ИВТб-2301-05-00 & \hspace{20mm}\rule[-1mm]{30mm}{0.10mm}\,/Черкасов А.А./ \\
                                 & \hspace{27.5mm}\footnotesize(подпись)                   \\
  Преподаватель                  & \hspace{20mm}\rule[-1mm]{30mm}{0.10mm}\,/Мельцов В. Ю./  \\
                                 & \hspace{27.5mm}\footnotesize(подпись)                   \\
\end{tabular}

\noindent
\begin{tabular}{lp{58mm}r}
  Работа защищена &  & \hspace{8mm}<<\rule[-1mm]{8mm}{0.10mm}\/>>\rule[-1mm]{38mm}{0.10mm}\ 2026 г.
\end{tabular}
\vfill

\begin{center}
  Киров\\
  2026
\end{center}

\newpage\thispagestyle{plain}

\section{Цель работы}
Получить практические навыки проектирования комбинационных логических устройств. Выполнить теоретический логический синтез и разработать принципиальную функциональную схему 4-разрядного цифрового компаратора, осуществляющего сравнение двух двоичных чисел.

\section{Условно-графическое обозначение (УГО)}
Цифровой компаратор принимает на вход два 4-разрядных слова $X (x_3, x_2, x_1, x_0)$ и $Y (y_3, y_2, y_1, y_0)$. На выходе формируются три взаимно исключающих сигнала: $R$ (Равно, $X=Y$), $B$ (Больше, $X>Y$) и $M$ (Меньше, $X<Y$).

\begin{figure}[H]
  \centering
  \begin{tikzpicture}[thick]
    \draw (0,-0.4) rectangle (3,4);
    \node at (1.5, 3.5) {$==$};
    % Входы X
    \draw (-0.5, 3.2) -- (0, 3.2) node[right] {$x_3$};
    \draw (-0.5, 2.8) -- (0, 2.8) node[right] {$x_2$};
    \draw (-0.5, 2.4) -- (0, 2.4) node[right] {$x_1$};
    \draw (-0.5, 2.0) -- (0, 2.0) node[right] {$x_0$};
    % Входы Y
    \draw (-0.5, 1.2) -- (0, 1.2) node[right] {$y_3$};
    \draw (-0.5, 0.8) -- (0, 0.8) node[right] {$y_2$};
    \draw (-0.5, 0.4) -- (0, 0.4) node[right] {$y_1$};
    \draw (-0.5, 0.0) -- (0, 0.0) node[right] {$y_0$};
    % Выходы
    \draw (3, 2.4) -- (3.5, 2.4) node[right] {$R$};
    \draw (3, 1.6) -- (3.5, 1.6) node[right] {$B$};
    \draw (3, 0.8) -- (3.5, 0.8) node[right] {$M$};
  \end{tikzpicture}
  \caption{УГО 4-разрядного цифрового компаратора}
\end{figure}

\section{Таблица истинности}
Полная таблица истинности устройства содержит $2^8 = 256$ строк. Для оптимизации описания использован метод фиксации безразличных состояний (обозначены знаком <<*>>). Сравнение чисел осуществляется поразрядно, начиная с самого старшего бита ($x_3$ и $y_3$).

\begin{center}
  \begin{longtable}{cccc|cccc|ccc}
    \caption*{Таблица 1 --- Таблица переходов 4-разрядного компаратора} \\
    \hline
    $x_3$ & $x_2$ & $x_1$ & $x_0$ & $y_3$ & $y_2$ & $y_1$ & $y_0$ & $R$ & $B$ & $M$ \\
    \hline
    \endfirsthead

    \caption*{(Продолжение таблицы 1)} \\
    \hline
    $x_3$ & $x_2$ & $x_1$ & $x_0$ & $y_3$ & $y_2$ & $y_1$ & $y_0$ & $R$ & $B$ & $M$ \\
    \hline
    \endhead

    \hline
    \endfoot

    \hline
    \endlastfoot

    1 & * & * & * & 0 & * & * & * & 0 & 1 & 0 \\
    0 & * & * & * & 1 & * & * & * & 0 & 0 & 1 \\
    0 & 1 & * & * & 0 & 0 & * & * & 0 & 1 & 0 \\
    0 & 0 & * & * & 0 & 1 & * & * & 0 & 0 & 1 \\
    1 & 1 & * & * & 1 & 0 & * & * & 0 & 1 & 0 \\
    1 & 0 & * & * & 1 & 1 & * & * & 0 & 0 & 1 \\
    0 & 0 & 1 & * & 0 & 0 & 0 & * & 0 & 1 & 0 \\
    0 & 0 & 0 & * & 0 & 0 & 1 & * & 0 & 0 & 1 \\
    1 & 1 & 1 & * & 1 & 1 & 0 & * & 0 & 1 & 0 \\
    1 & 1 & 0 & * & 1 & 1 & 1 & * & 0 & 0 & 1 \\
    1 & 0 & 1 & * & 1 & 0 & 0 & * & 0 & 1 & 0 \\
    1 & 0 & 0 & * & 1 & 0 & 1 & * & 0 & 0 & 1 \\
    0 & 1 & 1 & * & 0 & 1 & 0 & * & 0 & 1 & 0 \\
    0 & 1 & 0 & * & 0 & 1 & 1 & * & 0 & 0 & 1 \\
    0 & 0 & 0 & 1 & 0 & 0 & 0 & 0 & 0 & 1 & 0 \\
    0 & 0 & 0 & 0 & 0 & 0 & 0 & 1 & 0 & 0 & 1 \\
    1 & 1 & 1 & 1 & 1 & 1 & 1 & 0 & 0 & 1 & 0 \\
    1 & 1 & 1 & 0 & 1 & 1 & 1 & 1 & 0 & 0 & 1 \\
    0 & 0 & 1 & 1 & 0 & 0 & 1 & 0 & 0 & 1 & 0 \\
    0 & 0 & 1 & 0 & 0 & 0 & 1 & 1 & 0 & 0 & 1 \\
    0 & 1 & 0 & 1 & 0 & 1 & 0 & 0 & 0 & 1 & 0 \\
    0 & 1 & 0 & 0 & 0 & 1 & 0 & 1 & 0 & 0 & 1 \\
    1 & 0 & 0 & 1 & 1 & 0 & 0 & 0 & 0 & 1 & 0 \\
    1 & 0 & 0 & 0 & 1 & 0 & 0 & 1 & 0 & 0 & 1 \\
    0 & 1 & 1 & 1 & 0 & 1 & 1 & 0 & 0 & 1 & 0 \\
    0 & 1 & 1 & 0 & 0 & 1 & 1 & 1 & 0 & 0 & 1 \\
    1 & 0 & 1 & 1 & 1 & 0 & 1 & 0 & 0 & 1 & 0 \\
    1 & 0 & 1 & 0 & 1 & 0 & 1 & 1 & 0 & 0 & 1 \\
    1 & 1 & 0 & 1 & 1 & 1 & 0 & 0 & 0 & 1 & 0 \\
    1 & 1 & 0 & 0 & 1 & 1 & 0 & 1 & 0 & 0 & 1 \\
    0 & 0 & 0 & 0 & 0 & 0 & 0 & 0 & 1 & 0 & 0 \\
    1 & 1 & 1 & 1 & 1 & 1 & 1 & 1 & 1 & 0 & 0 \\
  \end{longtable}
\end{center}

\section{Канонические уравнения (КУ)}
На основе таблицы истинности составим систему логических уравнений для выходов $R$, $B$ и $M$. Выражение равенства одноимённых разрядов описывается эквивалентностью $(\overline{x}_i \cdot \overline{y}_i \vee x_i \cdot y_i)$.

Уравнение функции $R$ (Равно):
\begin{equation}
  R = (\overline{x}_3\overline{y}_3 \vee x_3y_3) \cdot (\overline{x}_2\overline{y}_2 \vee x_2y_2) \cdot (\overline{x}_1\overline{y}_1 \vee x_1y_1) \cdot (\overline{x}_0\overline{y}_0 \vee x_0y_0)
\end{equation}

Уравнение функции $B$ (Больше):
\begin{align*}
  B = {} & (x_3\overline{y}_3) \vee \\
         & (\overline{x}_3\overline{y}_3 \vee x_3y_3)(x_2\overline{y}_2) \vee \\
         & (\overline{x}_3\overline{y}_3 \vee x_3y_3)(\overline{x}_2\overline{y}_2 \vee x_2y_2)(x_1\overline{y}_1) \vee \\
         & (\overline{x}_3\overline{y}_3 \vee x_3y_3)(\overline{x}_2\overline{y}_2 \vee x_2y_2)(\overline{x}_1\overline{y}_1 \vee x_1y_1)(x_0\overline{y}_0)
\end{align*}

Уравнение функции $M$ (Меньше):
\begin{align*}
  M = {} & (\overline{x}_3y_3) \vee \\
         & (\overline{x}_3\overline{y}_3 \vee x_3y_3)(\overline{x}_2y_2) \vee \\
         & (\overline{x}_3\overline{y}_3 \vee x_3y_3)(\overline{x}_2\overline{y}_2 \vee x_2y_2)(\overline{x}_1y_1) \vee \\
         & (\overline{x}_3\overline{y}_3 \vee x_3y_3)(\overline{x}_2\overline{y}_2 \vee x_2y_2)(\overline{x}_1\overline{y}_1 \vee x_1y_1)(\overline{x}_0y_0)
\end{align*}

\section{Минимизация логических функций}
Анализ полученных канонических уравнений показывает, что функции не поддаются дальнейшей алгебраической минимизации. 

Функция $R$ не имеет повторяющихся термов. Для функций $B$ и $M$ невозможно вынести повторяющиеся блоки за скобки, так как условия проверки младших разрядов строго каскадно связаны с эквивалентностью всех предыдущих старших разрядов. Попытка сокращения приведёт к нарушению логики каскадной приоритетной проверки.

\section{Функциональная схема устройства}
Схема цифрового компаратора спроектирована на основе полученных неминимизированных канонических уравнений.

\begin{figure}[H]
  \centering
  \includegraphics[width=0.95\textwidth,keepaspectratio]{pics/functional_diagram.png}
  \caption{Схема электрическая функциональная 4-разрядного цифрового компаратора}
\end{figure}

\section*{Вывод}
В результате выполнения работы осуществлен структурный логический синтез 4-разрядного цифрового компаратора. На основе таблицы истинности с учетом безразличных состояний были выведены строгие канонические уравнения для функций <<Равно>>, <<Больше>> и <<Меньше>>. Доказано, что каскадный алгоритм поразрядного сравнения исключает возможность глубокой алгебраической минимизации из-за жесткой привязки младших разрядов к результатам проверки старших. Спроектированная функциональная схема полностью реализует заданный математический аппарат.
\end{document}